% 320_Spektraltheorie_En_ch.tex
% Detailed Field-Geometric Spectral Theory:
% Full Derivation of Particle Spectra, Mass Hierarchies
% and Gauge Groups from the T^4/Z_3 Geometry
% Doc. 320, August 2026
% Cross-references: Docs. 001, 003, 006, 007, 012, 133, 205, 317, 318, 319, 321, 322

% =========================================================
\section*{List of Symbols}
% =========================================================

\begin{center}
\begin{tabular}{lll}
\toprule
Symbol & Meaning & Value / Unit \\
\midrule
$\xi$        & FFGFT geometry parameter              & $4/30000 = 1.3\overline{3}\times10^{-4}$ \\
$v$          & Higgs vacuum expectation value        & $246.0$~GeV \\
$m_e$        & electron mass (neutrino sector ref.)  & $0.511$~MeV \\
$r_i$        & topological winding ratio             & $n_\phi^2/n_\theta$ \\
$p_i$        & exponent in mass formula              & dimensionless \\
$\eta_{D4}$  & $D_4$ projection coefficient          & $81/640$ \\
$F_{D4}$     & $D_4$ volume element                  & $27/1\,600\,000$ \\
$[K]$        & derivable from $\xi$                  & status marker \\
$[B]$        & algebraically proved                  & status marker \\
$[S]$        & sketched / plausible                  & status marker \\
\bottomrule
\end{tabular}
\end{center}

% =========================================================
\section*{Summary}
% =========================================================

The FFGFT derives all particle masses of the leptonic sector from
the single dimensionless geometry parameter $\xi = 4/30000$ and the
topology of the $T^4/\mathbb{Z}_3$ orbifold. This document follows
the derivation explicitly in every step: from the sub-Planck lattice,
through the $D_4$ projection and topological winding numbers, to the
measurable masses. The mathematical Hilbert-space foundation is provided
in Doc.~322.

\begin{keyresult}[Status markers]
All statements carry FFGFT status markers:\\
\textbf{[K]} derivable from $\xi$ \quad
\textbf{[B]} algebraically proved \quad
\textbf{[S]} sketched/plausible \quad
\textbf{[SETZUNG]} axiomatic setting
\end{keyresult}

% =========================================================
\section{Starting Point: the Sub-Planck Geometry}
% =========================================================

\subsection{Why a sub-Planck level?}

The Standard Model describes particle masses through Yukawa
couplings -- free parameters measured experimentally but not
explained within the framework itself. FFGFT goes one level deeper:
masses do not originate from free couplings but from the geometry of
space-time on scales below the Planck length $\ell_P$.

\textbf{[SETZUNG]} At the sub-Planck scale $t_0$ space-time has a
discrete packing structure. The most efficient regular lattice in four
dimensions is the $D_4$ lattice -- the 4D analogue of the densest
sphere packing in 3D.

\subsection{The $D_4$ lattice and its symmetry}

Properties of the $D_4$ lattice:
\begin{itemize}
  \item \textbf{4 orthogonal dimensions} with identical scaling.
  \item \textbf{Automorphism group} $\mathrm{Aut}(D_4) \cong W(F_4)$
    with $|W(F_4)| = 1152$ elements.
  \item \textbf{24 nearest neighbours} per lattice point.
  \item The torus $T^4/\mathbb{Z}_3$ is the natural compactification,
    with the $\mathbb{Z}_3$ identification reflecting the three quark
    colour charges.
\end{itemize}

\textbf{[K]} The fundamental FFGFT relation
$\tilde{T}\cdot m = 1$
follows from the duality between torus radius $R \propto \tilde{T}$
and mass $m \propto 1/R$: on a compact torus the smallest mode spacing
is inverse to the radius.

% =========================================================
\section{The Geometry Parameter $\xi$ and its Origin}
% =========================================================

\subsection{Definition and numerical value}

\textbf{[K]}
\begin{equation}
  \xi = \frac{4}{30000} = \frac{1}{7500}
      = 1.3\overline{3}\times10^{-4}.
  \label{eq:xi_def}
\end{equation}

\paragraph{Numerator 4} equals the number of orthogonal dimensions of
the $D_4$ lattice -- topologically fixed, not freely chosen.

\paragraph{Denominator 30000} equals the phase volume of the fractal
$T^4$ compactification: product of the $\mathbb{Z}_3$ orbifold order
($3$) and the number of fundamental Fourier modes on the torus
($10^4$, the fourth power of the first non-trivial mode).

\subsection{Connection to Higgs geometry}

\textbf{[K]} An independent derivation of $\xi$ follows from the Higgs
mass (Doc.~A130):
$\xi = \lambda_h^2 v^2 / (16\pi^3 m_h^2)$.
With the $K_{\mathrm{frac}}$-interface factor and iteration the fixed
point $\xi = 4/30000$ is obtained (Doc.~133). The Higgs path and the
$D_4$-lattice path are two independent routes to the same $\xi$.

\subsection{The \texorpdfstring{$D_4$}{D4} projection coefficient
  \texorpdfstring{$\eta_{D4}$}{eta\_D4}}

\textbf{[B]} The projection coefficient:
\begin{equation}
  \eta_{D4} = \frac{81}{640} = 0.1265625,
  \label{eq:eta_d4}
\end{equation}
and the $D_4$ volume element:
\begin{equation}
  F_{D4} = \xi\cdot\eta_{D4} = \frac{27}{1\,600\,000}
         = 1.6875\times10^{-5}.
  \label{eq:f_d4}
\end{equation}
Numerical check: $4\cdot81 = 324$, $\;30000\cdot640 = 19\,200\,000$,
$\;324/19\,200\,000 = 27/1\,600\,000$. \checkmark

% =========================================================
\section{Charged Leptons: Derivation from Winding Numbers}
% =========================================================

\subsection{Why winding numbers?}

Field modes on a compact torus $T^2$ are standing waves. The
periodicity condition $\psi(x + 2\pi R) = \psi(x)$ permits only
integer winding numbers $n \in \mathbb{Z}$. On the $T^2$ torus two
independent circles carry winding numbers $n_\theta$ and $n_\phi$.

\textbf{[K]} The three generations of charged leptons correspond to
the three topologically distinguished $(n_\theta, n_\phi)$ states
forming the simplest rational ratios $r_i = n_\phi^2/n_\theta$:

\begin{table}[h]
\centering
\caption{FFGFT quantum numbers of the charged leptons}
\label{tab:winding}
\begin{tabular}{lcccc}
\toprule
Particle & $n_\theta$ & $n_\phi$ & $r_i = n_\phi^2/n_\theta$ & $p_i$ \\
\midrule
Electron $e$  & 3 & 2 & $4/3$  & $3/2$ \\
Muon $\mu$    & 5 & 4 & $16/5$ & $1$   \\
Tau $\tau$    & 9 & 5 & $25/9$ & $2/3$ \\
\bottomrule
\end{tabular}
\end{table}

\subsection{The exponent ladder}

\textbf{[K]} The exponents $p_i$ form an arithmetic ladder with step
$\Delta p = -1/3$ per generation -- a direct imprint of the
$\mathbb{Z}_3$ orbifold order:
\begin{equation}
  p_e = \tfrac{3}{2},\quad p_\mu = 1,\quad p_\tau = \tfrac{2}{3}.
\end{equation}
The general mass formula:
\begin{equation}
  m_i = r_i\cdot\xi^{p_i}\cdot v,\qquad v = 246.0~\text{GeV}.
  \label{eq:mass_lepton}
\end{equation}

\subsection{Step-by-step mass calculation}

Pre-computed $\xi$ powers:
\begin{align}
  \xi^{3/2} &= 1.5396\times10^{-6}, &
  \xi^{1}   &= 1.3\overline{3}\times10^{-4}, &
  \xi^{2/3} &= 2.6099\times10^{-3}.
\end{align}

\textbf{[K]}
\begin{align}
  m_e   &= \tfrac{4}{3}\cdot1.5396\times10^{-6}\cdot246.0~\text{MeV}
         = \mathbf{0.511}~\mathbf{MeV}
         \quad(\text{PDG: }0.511~\text{MeV}),\\
  m_\mu &= \tfrac{16}{5}\cdot1.3\overline{3}\times10^{-4}
           \cdot246.0\times10^3~\text{MeV}
         = \mathbf{104.96}~\mathbf{MeV}
         \quad(\text{PDG: }105.66~\text{MeV}),\\
  m_\tau&= \tfrac{25}{9}\cdot2.6099\times10^{-3}
           \cdot246.0\times10^3~\text{MeV}
         = \mathbf{1783.5}~\mathbf{MeV}
         \quad(\text{PDG: }1776.86~\text{MeV}).
\end{align}

\subsection{Parameter-free mass ratios}

\textbf{[K]} The VEV $v$ cancels in ratios:
\begin{align}
  m_\mu/m_e   &= \tfrac{12}{5}\cdot\xi^{-1/2} = 207.85
               \quad(\text{PDG: }206.77,\;+0.52\,\%),\\
  m_\tau/m_\mu &= \tfrac{125}{144}\cdot\xi^{-1/3} = 16.99
               \quad(\text{PDG: }16.82,\;+1.04\,\%),\\
  m_\tau/m_e   &= \tfrac{25}{12}\cdot\xi^{-5/6} = 3532
               \quad(\text{PDG: }3477,\;+1.57\,\%).
\end{align}

\begin{table}[h]
\centering
\caption{Lepton masses and ratios: FFGFT vs.\ PDG}
\label{tab:leptons_comp}
\begin{tabular}{llllll}
\toprule
Quantity & Formula & Theory & PDG & Dev. & Status \\
\midrule
$m_e$          & $\frac{4}{3}\xi^{3/2}v$ & $0.511$~MeV   & $0.511$~MeV   & $<0.1\,\%$ & [K] \\
$m_\mu$        & $\frac{16}{5}\xi^1 v$   & $104.96$~MeV  & $105.66$~MeV  & $-0.66\,\%$ & [K] \\
$m_\tau$       & $\frac{25}{9}\xi^{2/3}v$ & $1783.5$~MeV & $1776.86$~MeV & $+0.38\,\%$ & [K] \\
$m_\mu/m_e$    & $\frac{12}{5}\xi^{-1/2}$ & $207.85$ & $206.77$ & $+0.52\,\%$ & [K] \\
$m_\tau/m_\mu$ & $\frac{125}{144}\xi^{-1/3}$ & $16.99$ & $16.82$ & $+1.04\,\%$ & [K] \\
\bottomrule
\end{tabular}
\end{table}

% =========================================================
\section{Neutrino Sector: Fixed-Point Localisation on the Orbifold}
% =========================================================

\subsection{Structural difference from charged leptons}

Charged leptons arise as winding modes on the $T^2$ torus and couple
to the Higgs VEV $v$. Neutrinos are fundamentally different: they are
localised at the \emph{geometric fixed points} of the
$T^4/\mathbb{Z}_3$ orbifold -- points $x_0$ satisfying
$\tau(x_0) = x_0$. This topological distinction produces a double
$\xi$-suppression compared to the charged leptons.
For the formal Hilbert-space treatment of the fixed-point boundary
conditions see Doc.~322, Eq.~(15).

\subsection{The neutrino mass formula}

\textbf{[K]}
\begin{equation}
  m_{\nu_i} = m_e\cdot\xi^{p_i},\qquad m_e = 0.511~\text{MeV}.
  \label{eq:nu_mass}
\end{equation}

\begin{table}[h]
\centering
\caption{Fixed-point localisation of neutrino flavours (Doc.~322)}
\label{tab:fixpoints}
\begin{tabular}{lllll}
\toprule
Fixed pt. & Position & Neutrino & Exponent $p_i$ & Mass \\
\midrule
$F_2$ & $\frac{2\pi R}{3}(1,1,1,0)$ & $\nu_1$ & $9/4$ & $0.976$~meV \\
$F_5$ & $\frac{4\pi R}{3}(1,1,1,0)$ & $\nu_2$ & $2$   & $9.084$~meV \\
$F_7$ & $\frac{4\pi R}{3}(1,1,1,1)$ & $\nu_3$ & $9/5$ & $54.11$~meV \\
\bottomrule
\end{tabular}
\end{table}

\subsection{Step-by-step neutrino mass calculation}

\textbf{[K]} Pre-computed $\xi$ powers and masses:
\begin{align}
  \xi^{9/4} &= 1.9103\times10^{-9},\quad
  m_{\nu_1} = 511\,000~\text{eV}\cdot1.9103\times10^{-9}
            = \mathbf{0.976}~\mathbf{meV},\\
  \xi^{2}   &= 1.\overline{7}\times10^{-8},\quad
  m_{\nu_2} = 511\,000~\text{eV}\cdot1.\overline{7}\times10^{-8}
            = \mathbf{9.084}~\mathbf{meV},\\
  \xi^{9/5} &= 1.059\times10^{-7},\quad
  m_{\nu_3} = 511\,000~\text{eV}\cdot1.059\times10^{-7}
            = \mathbf{54.11}~\mathbf{meV}.
\end{align}

\subsection{Mass-squared differences and oscillation dynamics}

\subsubsection{Solar mass-squared difference $\Delta m_{21}^2$}

\begin{equation}
  \Delta m_{21}^2 = (9.084~\text{meV})^2 - (0.976~\text{meV})^2
  = 82.52 - 0.95 = 81.57~\text{meV}^2
  = \mathbf{8.16\times10^{-5}~\text{eV}^2}.
\end{equation}
PDG~2023: $7.53\times10^{-5}$~eV$^2$. Deviation: $+8.3\,\%$.

\subsubsection{Atmospheric mass-squared difference $\Delta m_{32}^2$}

\begin{equation}
  \Delta m_{32}^2 = (54.11~\text{meV})^2 - (9.084~\text{meV})^2
  = 2928.3 - 82.5 = 2845.8~\text{meV}^2
  = \mathbf{2.846\times10^{-3}~\text{eV}^2}.
\end{equation}
PDG (NuFIT 5.3): $2.453\times10^{-3}$~eV$^2$. Deviation: $+16.0\,\%$.

\begin{tcolorbox}[colback=yellow!5!white, colframe=yellow!60!black,
  title={Open point: atmospheric mass-squared difference}]
$\Delta m_{21}^2$ agrees to $+8.3\,\%$, within the typical FFGFT range.
$\Delta m_{32}^2$ deviates by $+16.0\,\%$ -- about twice as much and with the
same sign: $m_{\nu_3}$ is too large. Equal sign and comparable magnitude
point to a common cause rather than two independent problems.

The fixed-point exponents $p_i$ are secured~[K]; the source of the
discrepancy is an open point. Three candidates were examined:
\begin{itemize}
  \item \textbf{$K_{\mathrm{frak}}$ correction at fixed point $F_7$}
    (Doc.~315): $F_7$ is the only fixed point in the full $(1,1,1,1)$
    direction and thus carries a mass-circle component. With
    $K_{\mathrm{frak}} = 1-100\xi = 74/75$ one obtains $m_3 = 53.39$~meV
    and $\Delta m_{32}^2 = 2.768\times10^{-3}$~eV², i.e.\ $+12.9\,\%$ --
    an improvement, but not sufficient.
  \item \textbf{Mixing term $F_5$--$F_7$}: the fixed points are not fully
    orthogonal. An off-diagonal term in the mass matrix on the orbifold
    would shift both eigenvalues and would carry the right sign. This
    requires the complete mass matrix, not worked out here -- the most
    promising candidate.
  \item \textbf{A different exponent} -- unlikely: the exponent required
    for the PDG value is $p_3 = 1.8081$, only $0.45\,\%$ above
    $9/5 = 1.8000$. No rational fraction with small numerators lies closer
    ($16/9 \to +74\,\%$, $11/6 \to -37.5\,\%$, $20/11 \to -17.1\,\%$).
    The exponent $9/5$ is the best rational candidate and is grounded in
    the $D_4$ coordination.
\end{itemize}
The deviation therefore most likely does not stem from a wrong exponent
but from an effect correcting $m_{\nu_i} = m_e\,\xi^{p_i}$
multiplicatively. Status: \textbf{[S]}.
\end{tcolorbox}

\subsubsection{Cosmological mass sum}

\begin{equation}
  \sum_i m_{\nu_i} = 0.976 + 9.084 + 54.11~\text{meV}
  = 64.17~\text{meV} = 0.0642~\text{eV}.
\end{equation}
Planck~2018 limit: $\sum m_\nu < 0.12$~eV. \checkmark

\begin{table}[h]
\centering
\caption{Neutrino spectrum: FFGFT prediction vs.\ experiment}
\label{tab:neutrinos}
\begin{tabular}{lllll}
\toprule
Quantity & Formula & Prediction & Experiment & Status \\
\midrule
$m_{\nu_1}$ & $m_e\xi^{9/4}$
  & $0.976$~meV & $<0.8$~eV  & [K] \\
$m_{\nu_2}$ & $m_e\xi^2$
  & $9.084$~meV & $\sim8.7$~meV & [K] \\
$m_{\nu_3}$ & $m_e\xi^{9/5}$
  & $54.11$~meV & $\sim50$~meV & [K] \\
$\Delta m_{21}^2$ & $m_{\nu_2}^2-m_{\nu_1}^2$
  & $8.16\times10^{-5}$~eV$^2$ & $7.53\times10^{-5}$~eV$^2$
  & [K] $+8.3\,\%$ \\
$\Delta m_{32}^2$ & $m_{\nu_3}^2-m_{\nu_2}^2$
  & $2.846\times10^{-3}$~eV$^2$ & $2.453\times10^{-3}$~eV$^2$
  & [S] $+16.0\,\%$ \\
$\sum m_\nu$ & $\sum m_{\nu_i}$
  & $0.0642$~eV & $<0.12$~eV & [K] \\
\bottomrule
\end{tabular}
\end{table}

% =========================================================
\section{Symmetry Breaking and Gauge Structure}
% =========================================================

\subsection{$SU(3)_c$: colour charges from $\mathbb{Z}_3$ fibration}

\textbf{[K]} The $\mathbb{Z}_3$ symmetry generates three discrete
complex-conjugate phase levels $\omega = e^{2\pi i/3}$, corresponding
to the three colour charges Red, Green, Blue. The eight generators of
$\mathfrak{su}(3)$ follow from the eight real degrees of freedom of
the $\mathbb{Z}_3$ projection. \textbf{[B]} Full algebraic derivation
of the Gell-Mann matrices from orbifold modes is in Doc.~321.

\subsection{Electroweak bosons ($W^\pm$, $Z^0$)}

\textbf{[B]}
\begin{equation}
  m_W = \tfrac{1}{2}\,g\,v \approx 80.38~\text{GeV},\qquad
  m_Z = m_W/\!\cos\theta_W \approx 91.19~\text{GeV}.
\end{equation}

% =========================================================
\section{Complete Overview of All Spectra}
% =========================================================

\begin{table}[h!]
\centering
\caption{Systematic overview of all FFGFT spectra, Doc.~320}
\label{tab:overview}
\begin{tabular}{lllll}
\toprule
\textbf{Group} & \textbf{Particle} & \textbf{Prediction} &
  \textbf{PDG} & \textbf{Status} \\
\midrule
Leptons
  & Electron $e$  & $0.511$~MeV   & $0.511$~MeV   & [K] \\
  & Muon $\mu$    & $104.96$~MeV  & $105.66$~MeV  & [K] \\
  & Tau $\tau$    & $1783.5$~MeV  & $1776.86$~MeV & [K] \\
\midrule
Neutrinos
  & $\nu_1$ ($F_2$) & $0.976$~meV  & $<0.8$~eV     & [K] \\
  & $\nu_2$ ($F_5$) & $9.084$~meV  & $\sim8.7$~meV & [K] \\
  & $\nu_3$ ($F_7$) & $54.11$~meV  & $\sim50$~meV  & [K] \\
\midrule
Oscillations
  & $\Delta m_{21}^2$
    & $8.16\times10^{-5}$~eV$^2$ & $7.53\times10^{-5}$~eV$^2$
    & [K] $+8.3\,\%$ \\
  & $\Delta m_{32}^2$
    & $2.846\times10^{-3}$~eV$^2$ & $2.453\times10^{-3}$~eV$^2$
    & [S] $+16.0\,\%$ \\
  & $\sum m_\nu$ & $0.0642$~eV & $<0.12$~eV & [K] \\
\midrule
Bosons
  & $\gamma$, $g$ & $0$ & $0$ & [K] \\
  & $W^\pm$       & $80.38$~GeV & $80.38$~GeV & [K] \\
  & $Z^0$         & $91.19$~GeV & $91.19$~GeV & [K] \\
\bottomrule
\end{tabular}
\end{table}

% =========================================================
\section{Cross-references and Open Bridges}
% =========================================================

\begin{table}[h]
\centering
\caption{Derivation chain Doc.~320 and cross-references}
\label{tab:refs}
\begin{tabular}{lll}
\toprule
Step & Doc. & Status \\
\midrule
$\xi = 4/30000$ from $D_4$ lattice            & Docs.~001, 205     & [K] \\
$\xi$ from Higgs geometry (second route)       & Doc.~A130          & [K] \\
$K_{\mathrm{frac}}$ interface factor           & Docs.~012, 133     & [K] \\
$\tilde{T}\cdot m = 1$                         & Doc.~003           & [SETZUNG] \\
Winding numbers $(n_\theta, n_\phi)$           & Doc.~317           & [K] \\
Lepton masses $m_e, m_\mu, m_\tau$             & Docs.~006, 317     & [K] \\
Hilbert-space embedding, operator $\hat{F}$    & Doc.~322           & [K/S] \\
Fixed-point positions $F_2, F_5, F_7$          & Doc.~322           & [K] \\
$\Delta m_{21}^2$ ($+8\,\%$)                  & this Doc.          & [K] \\
$\Delta m_{32}^2$ ($+16.0\,\%$)               & this Doc.          & [S] open \\
$SU(3)_c$ from $\mathbb{Z}_3$                  & Doc.~321           & [B] \\
Weinberg angle $\theta_W$ from $\xi$           & Doc.~323           & [K] \\
Quark sector (hadrons)                         & Docs.~318, 319     & open bridge \\
\bottomrule
\end{tabular}
\end{table}
